\documentclass[11pt,a4paper]{article}
\usepackage[margin=1in]{geometry}
\usepackage{amsmath,amssymb,amsthm}
\usepackage{algorithm}
\usepackage{algpseudocode}
\usepackage{booktabs}
\usepackage{hyperref}

\theoremstyle{definition}
\newtheorem{definition}{\textbf{Definition}}[section]
\newtheorem{theorem}{\textbf{Theorem}}[section]
\newtheorem{lemma}[theorem]{\textbf{Lemma}}
\newtheorem{remark}{\textbf{Remark}}[section]

\title{\textbf{Rainbow Extensions: Double DQN, Dueling Networks, and Prioritized Replay}}
\author{\textbf{Team Scaibu} \\ \small Email: \texttt{jaydeep.9664920749@gmail.com} \\ \small \textbf{Architecture and Core Engine Team}}
\date{\textbf{October 2026}}

\begin{document}
\maketitle

\section{Definitions and Cognitive Mental Models}

\subsection{Formal Mathematical Definition}
\textbf{Definition (Value-Advantage Factorization, Decoupled Action Selection, and Prioritized Sampling):}
Standard Deep Q-Networks suffer from three fundamental pathologies: overestimation bias of the maximization operator, redundant value representations across uninformative actions, and sample inefficiency from uniform experience replay. The \textbf{Rainbow Extensions} resolve these through three synergistic mathematical formulations:

\begin{enumerate}
    \item \textbf{Double Q-Learning (Decoupled Target Evaluation)}: Action selection is assigned to online parameters $\boldsymbol{\theta}$, while action evaluation is performed by target parameters $\boldsymbol{\theta}^-$:
    {\boldmath
    \begin{equation}
    y_t^{\text{Double}} = r_t + (1 - d_t) \gamma Q\left( s_{t+1}, \; \arg\max_{a' \in \mathcal{A}} Q(s_{t+1}, a'; \boldsymbol{\theta}); \; \boldsymbol{\theta}^- \right)
    \end{equation}
    }
    \item \textbf{Dueling Network Architecture}: The neural network decomposes state-action values into a scalar state-value stream $V(s) \in \mathbb{R}$ and an action-advantage stream $A(s, a) \in \mathbb{R}^{|\mathcal{A}|}$. To guarantee mathematical identifiability (so $V$ and $A$ cannot fluctuate arbitrarily by equal and opposite constants), Q-values are synthesized via mean-centering:
    {\boldmath
    \begin{equation}
    Q(s, a; \boldsymbol{\theta}, \alpha, \beta) = V(s; \boldsymbol{\theta}, \beta) + \left( A(s, a; \boldsymbol{\theta}, \alpha) - \frac{1}{|\mathcal{A}|}\sum_{a' \in \mathcal{A}} A(s, a'; \boldsymbol{\theta}, \alpha) \right)
    \end{equation}
    }
    \item \textbf{Prioritized Experience Replay (PER)}: Transitions are sampled non-uniformly proportional to their temporal difference surprise $p_t = (|\delta_t| + \epsilon)^\alpha$, with importance sampling weights $w_t = \left( N \cdot P(t) \right)^{-\beta} / \max_i w_i$ to correct sampling bias.
\end{enumerate}

\subsection{Expanded Non-Technical Definition (Intuition for Anyone)}
\textbf{Intuitive Conceptual Definition:}
\begin{enumerate}
    \item \textbf{Tackling the Three Flaws of Standard DQN}: Standard DQN was a major breakthrough, but it had three glaring weaknesses: it was overly optimistic (always picking the luckiest estimate), it wasted time calculating irrelevant details (trying to decide which button to press when falling off a cliff), and it reviewed boring game moments just as often as crucial winning plays.
    \item \textbf{Double DQN (Two Referees)}: Instead of letting one network pick the best play and grade itself, Double DQN has Referee 1 (online network) suggest the move, and Referee 2 (target network) evaluate its true worth. If Referee 1 got overly lucky on an action, Referee 2 keeps it grounded.
    \item \textbf{Dueling Networks (Knowing State Value Separately)}: In many driving games, whether you turn the steering wheel left or right on a long straight highway doesn't change your speed. Dueling networks calculate ``How good is this road?'' ($V$) once, rather than re-estimating the whole road for every steering angle ($A$).
    \item \textbf{Prioritized Replay (Studying What You Got Wrong)}: A smart student doesn't re-read easy flashcards they got right 100 times; they drill the 5 flashcards they failed on. PER replays surprise moments where the AI's predictions were wildly wrong until it masters them.
\end{enumerate}

\subsection{Child-Friendly Mental Model (Physical Metaphor)}
Imagine preparing for a chess tournament:
\begin{itemize}
    \item \textbf{The Two Grandmasters (Double DQN)}: You (Player 1) pick what you think is the best chess move. Your grandmaster coach (Player 2) inspects that exact move and grades it honestly so you don't overestimate your cleverness.
    \item \textbf{Board Advantage vs Move Choice (Dueling)}: Even before looking at specific moves, you recognize: ``I am up a Queen, so my overall board state is winning ($V$).'' Then you check individual moves ($A$).
    \item \textbf{The Bookmark Folder (PER)}: Every time you make a blunder in practice, you stick a bright yellow bookmark on that game board to review it ten times more often before tournament day!
\end{itemize}

\section{Core Mathematical Equations and Definitions}

\subsection{Decoupled Double Bellman Target}
{\boldmath
\begin{equation}
y^{\text{Double}} = r + \gamma \cdot Q\left( s', \; \arg\max_{a'} Q(s', a'; \boldsymbol{\theta}); \; \boldsymbol{\theta}^- \right)
\end{equation}
}
\begin{itemize}
    \item \textbf{Formal Definition}: The bootstrapped return target where argmax action selection is driven by online weights $\boldsymbol{\theta}$ and evaluation is performed by target weights $\boldsymbol{\theta}^-$.
    \item \textbf{What It Means}: Completely decouples the selection of the greedy action from the evaluation of its expected return.
    \item \textbf{Why It Is Important}: Eliminates upward maximization bias, yielding accurate, monotonic value convergence.
\end{itemize}

\subsection{Identifiable Dueling Q-Value Synthesis}
{\boldmath
\begin{equation}
Q(s, a) = V(s) + A(s, a) - \frac{1}{|\mathcal{A}|} \sum_{a' \in \mathcal{A}} A(s, a')
\end{equation}
}
\begin{itemize}
    \item \textbf{Formal Definition}: The mean-subtracted combination of the scalar state-value stream and action-advantage stream.
    \item \textbf{What It Means}: Forces the advantage stream to have zero mean across actions, making the state-value $V(s)$ identifiable as the average action value.
    \item \textbf{Why It Is Important}: Stabilizes optimization; updates to $V(s)$ immediately propagate to all action Q-values in a single gradient step.
\end{itemize}

\subsection{Prioritized Transition Sampling Distribution}
{\boldmath
\begin{equation}
P(i) = \frac{(|\delta_i| + \epsilon)^\alpha}{\sum_{k=1}^N (|\delta_k| + \epsilon)^\alpha}
\end{equation}
}
\begin{itemize}
    \item \textbf{Formal Definition}: The non-uniform categorical probability distribution over transition tuples in the replay buffer.
    \item \textbf{What It Means}: Samples transitions with large temporal difference errors more frequently than transitions that are already well-explained.
    \item \textbf{Why It Is Important}: Concentrates computational resources on the hardest learning examples, speeding up policy acquisition by $2\times$ to $5\times$.
\end{itemize}

\subsection{Importance Sampling Bias Correction Weight}
{\boldmath
\begin{equation}
w_i = \left( \frac{1}{N \cdot P(i)} \right)^\beta \cdot \frac{1}{\max_j w_j}
\end{equation}
}
\begin{itemize}
    \item \textbf{Formal Definition}: The normalized importance-sampling weight compensating for non-uniform transition sampling frequencies.
    \item \textbf{What It Means}: Down-weights frequently sampled high-error transitions during gradient updates to ensure unbiased expected gradient estimates.
    \item \textbf{Why It Is Important}: Prevents policy drift caused by over-sampling rare transition distributions.
\end{itemize}

\section{Rigorous Mathematical Analysis Checklist}

\begin{itemize}
    \item \textbf{Notation}: $V(s)$ (state value), $A(s, a)$ (action advantage), $Q(s, a)$ (action value), $\boldsymbol{\theta}$ (online network), $\boldsymbol{\theta}^-$ (target network), $\alpha$ (PER exponent), $\beta$ (PER annealing exponent).
    \item \textbf{Epistemic Status}: Proposed by Hessel et al. (DeepMind, AAAI 2018), integrates the 6 most impactful extensions of Deep Q-Networks.
    \item \textbf{Dependency Graph}: State $s \to [V(s), A(s, :)] \to Q(s, :) \to \text{Double DQN } y \to \delta \to \text{PER Priority } p$.
    \item \textbf{Dimensions}: Discrete action space $|\mathcal{A}|$, state value scalar $V(s) \in \mathbb{R}$, advantage vector $A(s, :) \in \mathbb{R}^{|\mathcal{A}|}$.
    \item \textbf{Regimes}: Value-based discrete control, complex gaming environments, combinatorial decision pipelines.
    \item \textbf{Invariants}: Identifiability constraint: $\frac{1}{|\mathcal{A}|}\sum_{a \in \mathcal{A}} (Q(s, a) - V(s)) = 0$.
    \item \textbf{Isomorphisms}: Dueling networks are isomorphic to an implicit baseline subtraction in policy gradients.
\end{itemize}

\section{Theoretical Theorems and Formal Proofs}

\begin{theorem}[Reduction of Overestimation Bias via Decoupled Double Q-Learning]
Let $Q_1(s', a) = Q^*(s', a) + \epsilon_{1, a}$ and $Q_2(s', a) = Q^*(s', a) + \epsilon_{2, a}$ be two independent estimators of true action-values $Q^*(s', a)$, where noise terms $\epsilon_{1, a}, \epsilon_{2, a} \sim \mathcal{N}(0, \sigma^2)$ are i.i.d. zero-mean random variables. Then the expected target value under Double Q-learning is strictly upper-bounded by standard Q-learning:
{\boldmath
\begin{equation}
\mathbb{E}\left[ Q_2\left( s', \; \arg\max_{a'} Q_1(s', a') \right) \right] \le \mathbb{E}\left[ \max_{a'} Q_1(s', a') \right]
\end{equation}
}
with the inequality holding strictly whenever $|\mathcal{A}| \ge 2$ and $\sigma^2 > 0$.
\end{theorem}

\begin{proof}
Let $a^* = \arg\max_{a'} Q_1(s', a')$.
In standard single Q-learning, the target evaluates:
{\boldmath
\begin{equation}
\mathbb{E}\left[ \max_{a'} Q_1(s', a') \right] = \mathbb{E}\left[ Q_1(s', a^*) \right]
\end{equation}
}
In Double Q-learning, the target evaluates $Q_2(s', a^*)$.
Taking the expectation:
{\boldmath
\begin{equation}
\mathbb{E}\left[ Q_2(s', a^*) \right] = \mathbb{E}\left[ \mathbb{E}\left[ Q_2(s', a^*) \;\middle|\; a^* \right] \right]
\end{equation}
}
Because $Q_2$ is independent of $Q_1$, the noise $\epsilon_{2, a^*}$ is independent of the action choice $a^*$ induced by $Q_1$.
Therefore:
{\boldmath
\begin{equation}
\mathbb{E}\left[ Q_2(s', a^*) \;\middle|\; a^* \right] = \mathbb{E}\left[ Q^*(s', a^*) + \epsilon_{2, a^*} \;\middle|\; a^* \right] = Q^*(s', a^*)
\end{equation}
}
Thus $\mathbb{E}[Q_2(s', a^*)] = \mathbb{E}[Q^*(s', a^*)]$.
In contrast, for standard Q-learning:
{\boldmath
\begin{equation}
\mathbb{E}\left[ Q_1(s', a^*) \right] = \mathbb{E}\left[ \max_{a'} \left( Q^*(s', a') + \epsilon_{1, a'} \right) \right] \ge \max_{a'} \mathbb{E}\left[ Q^*(s', a') + \epsilon_{1, a'} \right] = \max_{a'} Q^*(s', a')
\end{equation}
}
by Jensen's inequality (since the maximum function is convex).
When $|\mathcal{A}| \ge 2$ and $\sigma^2 > 0$, the maximum of multiple independent random variables strictly exceeds their individual expectations.
Therefore:
{\boldmath
\begin{equation}
\mathbb{E}\left[ Q_2\left( s', \; \arg\max_{a'} Q_1(s', a') \right) \right] \le \max_{a'} Q^*(s', a') < \mathbb{E}\left[ \max_{a'} Q_1(s', a') \right]
\end{equation}
}
proving that Double Q-learning strictly suppresses maximization overestimation bias.
\end{proof}

\begin{theorem}[Identifiability and Zero-Mean Centering in Dueling Networks]
Let $Q(s, a) = V(s) + A(s, a) - \frac{1}{|\mathcal{A}|}\sum_{a'} A(s, a')$. For any state $s$, the state-value $V(s)$ is uniquely identifiable as the arithmetic mean of action values across the action space:
{\boldmath
\begin{equation}
V(s) = \frac{1}{|\mathcal{A}|} \sum_{a \in \mathcal{A}} Q(s, a)
\end{equation}
}
\end{theorem}

\begin{proof}
Summing the synthesis equation over all actions $a \in \mathcal{A}$:
{\boldmath
\begin{align}
\sum_{a \in \mathcal{A}} Q(s, a) &= \sum_{a \in \mathcal{A}} \left[ V(s) + A(s, a) - \frac{1}{|\mathcal{A}|} \sum_{a' \in \mathcal{A}} A(s, a') \right] \\
&= \sum_{a \in \mathcal{A}} V(s) + \sum_{a \in \mathcal{A}} A(s, a) - \sum_{a \in \mathcal{A}} \left( \frac{1}{|\mathcal{A}|} \sum_{a' \in \mathcal{A}} A(s, a') \right) \\
&= |\mathcal{A}| \cdot V(s) + \sum_{a \in \mathcal{A}} A(s, a) - |\mathcal{A}| \cdot \frac{1}{|\mathcal{A}|} \sum_{a' \in \mathcal{A}} A(s, a') \\
&= |\mathcal{A}| \cdot V(s) + \sum_{a \in \mathcal{A}} A(s, a) - \sum_{a' \in \mathcal{A}} A(s, a') \\
&= |\mathcal{A}| \cdot V(s) + 0
\end{align}
}
Dividing both sides by $|\mathcal{A}| > 0$:
{\boldmath
\begin{equation}
V(s) = \frac{1}{|\mathcal{A}|} \sum_{a \in \mathcal{A}} Q(s, a)
\end{equation}
}
Thus $V(s)$ is uniquely defined as the average action value, resolving the parameter identifiability degeneracy.
\end{proof}

\section{Foundational Architectural Questions and Epistemic Meta-Questions}

\subsection{Architectural Question 1: Max-Centering vs Mean-Centering in Dueling Networks}
\textbf{Question:} Why does the Dueling architecture use mean-centering ($Q = V + A - \operatorname{mean} A$) instead of max-centering ($Q = V + A - \max A$)?\\
\textbf{Answer:} In theory, max-centering gives $V(s) = \max_a Q(s, a)$, directly matching the definition of the state value under an optimal policy. However, in practice, $\max A$ changes discrete argmax rapidly during training, making the gradient of the disadvantage discontinuous. Mean-centering ($\operatorname{mean} A$) provides smooth, stable gradients across all advantage streams simultaneously.

\textbf{Meta-Question:} Does mean-centering alter the relative ranking of actions?\\
\textbf{Answer:} No. Subtracting the constant $\frac{1}{|\mathcal{A}|}\sum A$ shifts all action values by the identical scalar offset, perfectly preserving the exact argmax action choice: $\arg\max_a Q(s, a) = \arg\max_a A(s, a)$.

\subsection{Architectural Question 2: Importance Sampling Annealing Schedule in PER}
\textbf{Question:} Why is the importance sampling exponent $\beta$ annealed from $\beta_0 \approx 0.4$ to $\beta_{\text{final}} = 1.0$ over training?\\
\textbf{Answer:} Early in training, the prioritized transitions provide rich diverse gradient signals, and small estimation bias does not hinder rapid exploratory learning. Near convergence, exact unbiased gradient estimation is mandatory to ensure the value function converges precisely to the true Bellman fixed point $Q^*$, requiring $\beta = 1.0$ for full importance sampling correction.

\textbf{Meta-Question:} What are the other components of the full Rainbow framework?\\
\textbf{Answer:} Full Rainbow integrates six components: Double DQN, Dueling Networks, Prioritized Replay, Multi-step Bootstrap ($n$-step returns), Distributional Value Learning (C51 categorical distribution), and Noisy Networks (parametric exploration).

\section{Algorithmic Implementation Specification}

\begin{algorithm}[H]
\caption{Rainbow Dueling Synthesis, Double DQN Target, and PER Priority}
\begin{algorithmic}[1]
\Require State value $V \in \mathbb{R}$, advantages $\mathbf{A} \in \mathbb{R}^{|A|}$, action $a$, reward $r$, online next Q $\mathbf{q}_{\text{on}} \in \mathbb{R}^{|A|}$, target next Q $\mathbf{q}_{\text{tar}} \in \mathbb{R}^{|A|}$, discount $\gamma$, terminal flag $done$, priority exponent $\alpha$
\Ensure Synthesized Q-values $\mathbf{Q} \in \mathbb{R}^{|A|}$, Double target $y_{\text{Double}}$, TD error $\delta$, priority $p$, greedy action $a^*$
\State $\mu_A \gets \frac{1}{|A|}\sum_{i=0}^{|A|-1} \mathbf{A}[i]$ \Comment{Dueling Mean-Centering}
\State Initialize $\mathbf{Q} \in \mathbb{R}^{|A|}$
\For{$i = 0$ \textbf{to} $|A|-1$}
    \State $\mathbf{Q}[i] \gets V + (\mathbf{A}[i] - \mu_A)$
\EndFor
\If{$done$ is true} \Comment{Double DQN Target Evaluation}
    \State $y_{\text{Double}} \gets r$
\Else
    \State $a_{\text{best}} \gets \arg\max_{0 \le i < |A|} \mathbf{q}_{\text{on}}[i]$ \Comment{Action Selected by Online Network}
    \State $y_{\text{Double}} \gets r + \gamma \cdot \mathbf{q}_{\text{tar}}[a_{\text{best}}]$ \Comment{Value Evaluated by Target Network}
\EndIf
\State $pred \gets \mathbf{Q}[a]$
\State $\delta \gets y_{\text{Double}} - pred$ \Comment{Double TD Error}
\State $p \gets (|\delta| + 10^{-6})^\alpha$ \Comment{Prioritized Experience Replay Priority}
\State $a^* \gets \arg\max_{0 \le i < |A|} \mathbf{Q}[i]$
\State \Return $(\mathbf{Q}, y_{\text{Double}}, \delta, p, a^*)$
\end{algorithmic}
\end{algorithm}

\section{Big-$\Theta$ Complexity Analysis}

\begin{itemize}
    \item \textbf{Dueling Synthesis Time Complexity}: $\Theta(|A|)$ scalar operations to average advantages and compute $\mathbf{Q}$.
    \item \textbf{Double Target Evaluation Complexity}: $\Theta(|A|)$ operations to find the online maximum action.
    \item \textbf{Priority Calculation}: $\Theta(1)$ operations.
    \item \textbf{Space Complexity}: $\Theta(|A|)$ working buffer for synthesized Q-values.
    \item \textbf{Parallel Depth}: $\mathcal{D} = \Theta(\log |A|)$.
\end{itemize}

\section{Single Canonical Repository Reference}

The complete deterministic scalar implementation, verification suite, and unit test benchmarks are published at:
\begin{center}
\url{https://github.com/Chief-Strategist-J/llm-obs-infra}
\end{center}

\end{document}
